% ===================================================================%
% ===================================================================
\section{\SHACL\ as a Unified Data Modelling Language}
% ===================================================================

\subsection{The Dual Role of \SHACL}

A central innovation of the proposed architecture is the use of \SHACL\ (Shapes
Constraint Language)~\cite{shacl_spec} not merely as a validation tool but as a
\textbf{unified data modelling language} that serves two complementary purposes:
(1)~as the basis for conformity checking of submitted \DPP\ data against
applicable regulations and standards, and (2)~as the generative source for user
interfaces --- \DPP\ websites and application program interfaces --- via
\SHACL-to-HTML rendering engines. This dual role transforms \SHACL\ from a
passive validation layer into an active modelling substrate that drives both
data ingestion and data presentation.

At \DPPforEU~2026, \SHACL\ was described by a conference participant as
``\textit{a very elegant approach}'' for \DPP\ validation, and a NILU researcher
was introduced as speaking about ``\textit{\SHACL\ information models and how
they will be solving all our problems}''~\cite{dpp4eu2026_day1}. The conference
discussion revealed that the community recognises \SHACL's potential but has not
yet connected it to a portal-backed architecture where \SHACL\ shapes serve as
the single source of truth for both validation and interface generation. The
proposed architecture makes this connection explicit: \SHACL\ shapes are stored
in the public repository on the \EC\ \DPP\ registry, used by the portal
controller for ingestion validation (Section~4.2.1), and used by application
developers to generate user interfaces that, by construction, ensure
user-entered data conforms to regulations and standards.

\subsection{\SHACL\ for Conformity Checking}

The first role of \SHACL\ is conformity checking: verifying that submitted \DPP\
data satisfies the structural, cardinality, datatype, and semantic constraints
defined by applicable regulations and harmonised standards. As a \JTC~24
representative observed at \DPPforEU~2026, machine-readable product data must
become checkable by machine-readable regulation~\cite{dpp4eu2026_day1}. The
portal-backed architecture makes this a present capability by maintaining
\SHACL\ shapes --- one set for each product category and applicable regulation
--- in the same semantic framework as the \DPP\ data itself.

When an \REO\ submits \DPP\ data to the portal, the Ingestion Governance Loop
validates the data against the \SHACL\ shapes corresponding to the product's
category and regulatory context (Section~4.2.1). If validation fails, the \DPP\
is rejected with a structured error report indicating which shapes were violated
and which data points failed. This design ensures that only regulation-compliant
data enters the triple-store. The conformity checking is not a post-hoc audit
but a precondition for ingestion: data that does not conform to the applicable
\SHACL\ shapes is never committed to the portal.

The conference observation that ``\textit{you still have to look in the book
whether what you see on the screen is actually according to
requirements}''~\cite{dpp4eu2026_day1} is resolved structurally: the ``book''
(the regulation) is represented as \SHACL\ shapes, and the checking is performed
automatically at ingestion time. The quality infrastructure community's emphasis
on linking \DPPs\ to standards, measurements, and certificates (see Section~7.4)
is realised through \SHACL\ shapes that reference the applicable standards and
certificate schemas, ensuring that the \DPP\ data is not only structurally valid
but also semantically aligned with the regulatory requirements it claims to
satisfy.

\subsection{\SHACL\ for \GUI\ Generation via \DASH}

The second role of \SHACL\ is interface generation: using \SHACL\ shapes as the
basis for generating user interfaces --- web forms, mobile app screens, and
\API\ schemas --- that enable economic operators to enter \DPP\ data. This is
implemented through \SHACL-to-HTML rendering engines such as \DASH\ (\SHACL\
Application Dashboard)~\cite{dash_shacl} and the open-source shacl-form
library~\cite{shacl_form_lib} developed by the Universitäts- und
Landesbibliothek Darmstadt. At \DPPforEU~2026, a participant from the DigiPass
project described how \SHACL\ can ``\textit{validate a knowledge graph that is
coming as ingested information and it can also be used as has been done in our
project to automatically generate a form for manual data
ingest}''~\cite{dpp4eu2026_day2}.

This dual-use capability is the key insight: the same \SHACL\ shapes that
validate submitted data at the portal can also generate the data entry forms
that produce the data in the first place. When an application developer builds a
\DPP\ data entry tool --- a web platform or a mobile app --- they use the
\SHACL\ shapes from the public repository to generate the user interface. The
generated interface presents fields corresponding to the \SHACL\ properties,
enforces cardinality and datatype constraints, and validates user input against
the shapes in real time. By construction, the data entered through such an
interface conforms to the regulations and standards encoded in the \SHACL\
shapes. This eliminates the gap between data entry and data validation: the same
shapes govern both, ensuring that data is ``born compliant'' rather than being
checked for compliance after the fact.

\subsection{The \SHACL\ Brick Library}

To enable modular, maintainable, and regulation-aware interface generation, we
propose a \textbf{\SHACL\ brick library}: a collection of modular,
regulation-compliant \SHACL\ shape components --- ``bricks'' --- that are
reviewed against applicable regulations, delegated acts, and harmonised
standards, and then assembled into complete \SHACL\ schemas for specific product
categories and regulatory contexts.

A \SHACL\ brick represents a coherent unit of regulatory or standard-based
requirements: for example, a ``chemical composition brick'' that encodes the
data points and constraints required by \REACH\ for substance declaration, a
``carbon footprint brick'' that encodes the data points required by the ESPR for
environmental impact declaration, or a ``battery composition brick'' that
encodes the data points required by the Battery Regulation. Each brick is
independently developed, independently validated against the source regulation,
and independently versioned. When a new regulation is adopted, or an existing
regulation is amended, only the affected bricks are updated; all schemas that do
not use those bricks remain unchanged.

The brick library is hosted on the public repository for ontologised regulations
and standards (Section~6.7), making it accessible to all \DPP\ stakeholders:
portal operators, application developers, economic operators, and market
surveillance authorities. The conference principle that ``\textit{everything
needs to be ontologised}'' (Section~3.5) is realised through the brick library:
regulations, standards, data models, and validation shapes are all represented
as modular semantic artefacts within a unified framework.

\subsection{The Two-Stage \SHACL\ Design Toolchain}

The generation of complete \SHACL\ schemas from the brick library requires a
\textbf{two-stage \SHACL\ design toolchain}: a brick generator and a schema
assembler.

\subsubsection{Stage 1: The Brick Generator}

The brick generator is a tool that takes ontologised regulatory and
standard-based information as input and produces individual \SHACL\ bricks as
output. The input is domain-specific semantic information: ontologies,
vocabularies, and structured representations of regulatory requirements for a
specific product domain (e.g., batteries, textiles, construction products,
chemicals). The output is a set of \SHACL\ shape components, each encoding a
coherent unit of regulatory requirements. The brick generator requires that
semantic information for product-domain-specific regulations and standards be
generated and made available (Section~6.8). This is a non-trivial task: it
requires domain experts (regulatory specialists, standards experts, product
category specialists) to collaborate with ontology engineers to produce
machine-readable representations of regulatory requirements. The conference
discussion on building ``\textit{ontologies with
experts}''~\cite{dpp4eu2026_day2} directly motivates this stage: the brick
generator is the tool that bridges the gap between domain expert knowledge and
machine-readable \SHACL\ shapes.

\subsubsection{Stage 2: The Schema Assembler}

The schema assembler is a tool that takes a set of \SHACL\ bricks and assembles
them into a complete \SHACL\ schema for a specific product category and
regulatory context. The assembler selects the relevant bricks based on the
product category, the applicable regulations, and the applicable harmonised
standards, and combines them into a unified \SHACL\ shape graph. The assembled
schema is then used for two purposes: (1)~by the portal controller for ingestion
validation, and (2)~by application developers for interface generation via
\DASH\ or shacl-form. The assembler ensures that the assembled schema is
consistent --- no conflicting constraints, no duplicate properties, no missing
mandatory shapes --- and that it covers all regulatory requirements applicable
to the target product category. When a regulation is updated, the assembler
re-generates the schema from the updated bricks, ensuring that all downstream
applications --- portal validation, data entry forms, \API\ schemas --- reflect
the most current regulatory requirements.

\subsection{Licensed \DPP\ Application Ecosystem}
\label{sec:licensed-apps}

The \SHACL-driven architecture creates a new category of economic actor:
\textbf{licensed \DPP\ application providers} --- software businesses that
specialise in building, maintaining, and distributing regulation-compliant \DPP\
generation and maintenance applications for specific product domains. Rather
than monolithic \DPP\ platforms, these are lightweight, domain-specific
applications generated directly from the regulation-conformant \SHACL\ schemas
assembled by the schema assembler. For large
companies, the \DPP\ task may be handled as part of their internal
digitalisation programme; for smaller businesses, however, generating a \DPP\
can be very costly in terms of the time required to learn about the regulations
and delegated acts, and the process of registering and
maintenance~\cite{dpp4eu2026_postconf}. Licensed \DPP\ application providers
fill this gap by offering expertly designed software tools that run on small
computer systems, including PCs and mobile phones, handling the full \DPP\
lifecycle --- creation, registration, update, and maintenance --- without
requiring the economic operator to develop deep regulatory or semantic
expertise. As the DigiPass consortium envisioned, the approved \SHACL\ snippet
library serves as a very valuable source for generating such
applications~\cite{dpp4eu2026_postconf}.

\subsubsection{The Three-Stage Licensing Pipeline}

The generation of licensed \DPP\ applications extends the two-stage \SHACL\
design toolchain (Section~6.5) with a third stage: \textbf{app generation,
review, and licensing}. The full pipeline is:

\begin{enumerate}
\item \textbf{Brick generation} (Stage~1): domain experts and ontology engineers
produce \SHACL\ bricks from ontologised regulations and domain-specific semantic
information.
\item \textbf{Schema assembly} (Stage~2): the schema assembler combines bricks
into complete \SHACL\ schemas for specific product categories and regulatory
contexts.
\item \textbf{App generation and licensing} (Stage~3): licensed \DPP\
application providers construct tailored applications from the assembled \SHACL\
schemas, using \SHACL-to-HTML rendering engines such as \DASH\ or shacl-form.
The resulting application is reviewed, approved, and licensed for distribution.
\end{enumerate}

\subsubsection{Dual Authorisation}

The licensing of \DPP\ application providers follows a dual authorisation model,
reflecting the dual nature of the regulatory and professional requirements.
First, the \EC, through the portal authority, authorises the software business
to operate as a licensed \DPP\ application provider, verifying its legal
existence and operational competence. Second, professional organisations (e.g.,
industry associations, standardisation bodies) authorise the provider's
domain-specific competence, ensuring that the applications it produces are
technically and regulatorily appropriate for the product categories they cover.
This dual authorisation model mirrors the trust tier matrix (Section~5.5): legal
existence is verified through the first pillar, and domain competence is
verified through the second pillar, but applied to software providers rather
than data requesters.

\subsubsection{The Review Board}

A \textbf{review board} --- a technical body established at the authority/portal
level --- validates submitted \SHACL\ bricks, assembled schemas, and generated
applications against the applicable regulations, delegated acts, and
norms~\cite{dpp4eu2026_postconf}. The review board has three responsibilities:

\begin{enumerate}
\item \textbf{Brick review}: each \SHACL\ brick submitted to the public library
is reviewed for regulatory accuracy, semantic correctness, and consistency with
the source regulation.
\item \textbf{Schema review}: assembled \SHACL\ schemas are reviewed for
completeness (all applicable regulatory requirements covered), consistency (no
conflicting constraints), and correctness (no missing mandatory shapes).
\item \textbf{Application review}: generated \DPP\ applications are reviewed for
regulatory compliance, correct use of the portal \API, and foolproofness (see
Section~6.6.4).
\end{enumerate}

The review board is funded through a dedicated \EU\ budget line that finances
the maintenance of the open-access standards repository, the \SHACL\ brick
review process, and the federation-service \APIs~\cite{dpp4eu2026_postconf}.

\subsubsection{Foolproof by Construction}

The defining property of a licensed \DPP\ application is that it is
\textbf{foolproof by construction}: the application generates \RDF\ graphs that,
when submitted to the portal, will not fail \SHACL\ validation. This guarantee
is structural, not accidental. Because the application's data entry forms are
generated from the same \SHACL\ shapes that the portal controller uses for
ingestion validation (Section~4.2.1), the forms enforce exactly the constraints
that the portal will check. Cardinality, datatype, and semantic constraints are
validated in real time during data entry; the user cannot submit a \DPP\ that
violates a shape constraint because the form itself prevents it. As a conference
participant from the DigiPass project described, \SHACL\ can ``\textit{validate
a knowledge graph that is coming as ingested information and it can also be used
as has been done in our project to automatically generate a form for manual data
ingest}''~\cite{dpp4eu2026_day2}. The licensed \DPP\ application ecosystem
operationalises this dual-use capability as a guarantee: the same shapes that
validate at the portal generate the forms at the point of data entry, ensuring
that data is ``born compliant'' (Section~6.3). This is why the architecture
favours lightweight, domain-specific apps over monolithic platforms: only
applications generated from the same \SHACL\ shapes that govern portal ingestion
can offer a structural guarantee of compliance.

\subsubsection{Portal \API\ Integration}

Licensed \DPP\ applications must integrate with the portal \API\ (Section~4.6)
for all \DPP\ lifecycle operations: authentication, \DPP\ creation, update,
revocation, registration with the \EC\ central registry, link traversal, and
\SPARQL\ query. This requirement ensures that:

\begin{itemize}
\item All \DPP\ data generated by licensed applications passes through the
Ingestion Governance Loop, receiving \SHACL\ validation, cryptographic
alignment, and external entity normalisation.
\item All \DPP\ updates propagate correctly through the linked \DPP\ network via
\texttt{dpp:hasPrecursor} and \texttt{dpp:derivedFrom}.
\item All access requests pass through the Resolution Governance Loop, with
\ABAC\ policies enforced uniformly regardless of the originating application.
\end{itemize}

The portal \API\ thus serves as the single integration point between licensed
applications and the portal infrastructure, ensuring architectural consistency
across the entire application ecosystem.

\subsection{The Public Repository for Ontologised Regulations and Standards}

The brick library and the assembled \SHACL\ schemas are hosted in a
\textbf{public repository for ontologised regulations and standards}, maintained
on the \EC\ \DPP\ registry alongside the central identifier resolution service.
This repository contains:

\begin{itemize}
\item \textbf{Ontologised regulations}: machine-readable representations of the
ESPR, the Battery Regulation, \REACH, and product-specific delegated acts,
expressed as \RDF\ graphs with links to the original legal text.
\item \textbf{Ontologised standards}: machine-readable representations of
harmonised standards (EN~18215--18232 \cite{en18215}) and sector-specific
standards, expressed as \RDF\ graphs.
\item \textbf{\SHACL\ bricks}: modular shape components encoding regulatory and
standard-based requirements, produced by the brick generator.
\item \textbf{Assembled \SHACL\ schemas}: complete shape graphs for specific
product categories and regulatory contexts, produced by the schema assembler.
\item \textbf{Domain ontologies}: product-domain-specific ontologies (e.g.,
battery ontology, textile ontology, construction ontology) that provide the
semantic vocabulary for the \SHACL\ shapes.
\end{itemize}

The repository is public, freely accessible, and continuously updated as
regulations evolve and new product categories come under \DPP\ requirements. The
conference aspiration that standards should be ``\textit{publicly accessible,
and they should be for free, they should not be sold by \DIN\ and other
organisations}''~\cite{dpp4eu2026_day2} is operationalised through this
repository: ontologised regulations and standards are hosted on a public \EC\
infrastructure, accessible to all \DPP\ stakeholders without payment. Both the
regulations (as ontologised legal text) and the validation shapes (as \SHACL\
bricks and schemas) are available in machine-readable form, enabling automated
compliance checking at scale.

\subsection{Domain-Specific Semantic Information for Brick Generation}

The brick generator (Stage~1 of the toolchain) requires that semantic
information for product-domain-specific regulations and standards be generated
and made available. This is the foundational layer of the \SHACL\ architecture:
without domain-specific ontologies and structured regulatory representations,
the brick generator has no input, and no \SHACL\ bricks can be produced. As the
Materials Commons presentation at \DPPforEU~2026 described, this involves
``\textit{the implementation of the semantic framework network and the
infrastructure setup}''~\cite{dpp4eu2026_day2}.

The domain-specific semantic information must cover:

\begin{itemize}
\item \textbf{Regulatory requirements}: the specific data points, thresholds,
and constraints mandated by the ESPR, delegated acts, and product-specific
regulations for each product category.
\item \textbf{Standard-based requirements}: the specific data points,
measurement methods, and conformity assessment procedures defined by harmonised
standards (EN~18215--18232) \cite{en18215, en18216, en18217, en18221, en18222}
and sector-specific standards.
\item \textbf{Data models}: the structured representations of product,
component, and material data for each product category, including properties,
relationships, and cardinalities.
\item \textbf{Vocabularies}: the controlled vocabularies, code lists, and
classification schemes used in each product domain (e.g., material
classification codes, hazard categories, recycling codes).
\end{itemize}

This domain-specific semantic information is the input to the brick generator,
which produces \SHACL\ bricks that encode these requirements in a
machine-actionable form. The bricks are then assembled into complete schemas by
the schema assembler and published in the public repository, where they are
available for portal validation and interface generation.

\subsection{The SHACL-Driven Data Lifecycle}

The \SHACL-centric architecture creates a coherent data lifecycle that spans
from regulation to data entry to validation to storage to access:

\begin{enumerate}
\item \textbf{Regulation ontologisation}: regulations and standards are
represented as machine-readable semantic artefacts in the public repository.
\item \textbf{Brick generation}: the brick generator produces \SHACL\ bricks
from the ontologised regulations and domain-specific semantic information.
\item \textbf{Schema assembly}: the schema assembler combines bricks into
complete \SHACL\ schemas for specific product categories.
\item \textbf{Application generation and licensing}: licensed \DPP\ application
providers construct tailored applications from the assembled schemas, reviewed
and approved by the review board (Section~6.6).
\item \textbf{Data entry}: economic operators enter \DPP\ data through the
licensed applications, which validate input against the \SHACL\ shapes in real
time.
\item \textbf{Ingestion validation}: the portal controller validates submitted
\DPP\ data against the same \SHACL\ shapes during the Ingestion Governance Loop.
\item \textbf{Storage}: validated \DPP\ data is committed to the portal's
triple-store as \RDF\ named graphs.
\item \textbf{Access}: the Resolution Governance Loop filters and serves \DPP\
data based on \ABAC\ policies (Section~5).
\end{enumerate}

At every stage, the \SHACL\ shapes serve as the single source of truth for what
constitutes valid, regulation-compliant \DPP\ data, making the conference
aspirations for machine-checkable regulation and universal ontologisation
operational within a single, coherent architectural framework.
